\chapter{Salpingo-Oophorectomy}

\section*{Introduction \& Clinical Indications}
A salpingo-oophorectomy is a surgical procedure that involves the excision of one or both fallopian tubes (salpinx) and ovaries (oophoron). This operation is primarily performed within the female pelvis, targeting the adnexal structures, which are critical for female reproductive health. Salpingo-oophorectomy is indicated in various clinical scenarios, including benign conditions such as large ovarian cysts, endometriosis, and tubo-ovarian abscesses, as well as malignant diseases like ovarian or fallopian tube cancer. The clinical significance of this procedure extends beyond treatment; it plays a vital role in risk reduction for individuals with a genetic predisposition to certain cancers, particularly those with BRCA1 or BRCA2 mutations. By removing these structures, the surgery not only addresses existing pathology but also prevents future disease, with bilateral procedures inducing surgical menopause and its associated hormonal implications.

\begin{itemize}
  \item Unilateral Salpingo-Oophorectomy (USO)
  \item Bilateral Salpingo-Oophorectomy (BSO)
  \item Adnexectomy
  \item Ovarian and Tubal Excision
  \item Risk-Reducing Salpingo-Oophorectomy (RRSO)
  \item Oophorectomy with Salpingectomy
\end{itemize}

The fundamental principle of salpingo-oophorectomy is the complete surgical excision of the affected fallopian tube and ovary from their anatomical attachments. This involves careful identification and ligation of the vascular pedicles, particularly the infundibulopelvic ligament that contains the ovarian artery and vein, as well as the utero-ovarian ligament. 

The rationale for performing this surgery varies based on the underlying indication. In benign conditions, the primary goal is to alleviate symptoms, prevent complications such as torsion or rupture, and remove any inflammatory or neoplastic processes. In cases of malignancy, the procedure aims to excise the primary tumor, facilitate accurate cancer staging, and contribute to cytoreduction efforts.

For risk-reducing surgeries, the objective is to eliminate the target organs where ovarian and fallopian tube cancers typically arise, significantly lowering the lifetime risk of developing these malignancies. Technical goals include achieving complete removal of the adnexa, ensuring hemostasis to minimize blood loss, and meticulously avoiding injury to adjacent vital structures such as the ureters, bowel, and major pelvic vessels. The decision to perform unilateral versus bilateral removal, as well as the choice of surgical approach (laparoscopic, robotic, or open), is guided by the specific pathology, the patient's age, fertility desires, and overall medical status.

The history of salpingo-oophorectomy dates back to the early 19th century, with Dr. Ephraim McDowell performing the first successful oophorectomy in 1809 in Kentucky for a massive ovarian tumor. Initially, this procedure was fraught with risks due to the absence of anesthesia and antiseptic techniques, making it a dangerous undertaking. As surgical techniques and anatomical knowledge improved, oophorectomy became more common in the late 19th century, primarily performed via open laparotomy.

The late 20th century marked a significant evolution in the field with the introduction of minimally invasive surgical techniques. Dr. Kurt Semm pioneered laparoscopic approaches in gynecology during the 1970s, leading to the widespread adoption of laparoscopic salpingo-oophorectomy by the 1990s. Concurrently, advancements in genetic research, particularly regarding BRCA1 and BRCA2 mutations, facilitated the development of risk-reducing salpingo-oophorectomy as a prophylactic measure in high-risk populations, further refining the indications and techniques for this essential gynecologic surgery.

Salpingo-oophorectomy is indicated in various clinical scenarios, including:

\begin{itemize}
  \item Removal of persistent, symptomatic, or suspicious ovarian cysts or masses that do not resolve spontaneously or raise concerns for malignancy.
  \item Definitive treatment of benign ovarian tumors, such as endometriomas, dermoid cysts (teratomas), or fibromas, that cause significant pain or mass effect.
  \item Surgical intervention for ovarian torsion, a condition where the ovary twists on its vascular pedicle, compromising blood supply and requiring urgent treatment.
  \item Management of severe endometriosis involving the adnexa, particularly when medical management fails or large endometriomas are present.
  \item Treatment of tubo-ovarian abscess (TOA) unresponsive to prolonged antibiotic therapy or in cases of rupture.
  \item Primary surgical treatment for confirmed or suspected ovarian, fallopian tube, or primary peritoneal cancer, often as part of a staging or debulking procedure.
  \item Prophylactic or risk-reducing surgery in individuals with a high genetic risk for ovarian or fallopian tube cancer, such as those with BRCA1, BRCA2, or Lynch syndrome mutations.
  \item Prevention of future adnexal pathology, sometimes performed concurrently with a hysterectomy for other indications in perimenopausal or postmenopausal women.
  \item Rarely, for extensive damage from an ectopic pregnancy that cannot be managed by salpingectomy alone.
\end{itemize}

Salpingo-oophorectomy is selected for specific adnexal conditions, such as torsion with a nonviable or severely damaged adnexa, selected symptomatic masses, suspected or confirmed malignancy, and risk-reducing surgery in people at high genetic risk. Ovarian conservation or cystectomy may be preferable in younger patients with benign disease when feasible.

However, salpingo-oophorectomy can also function as a secondary or adjunct procedure. It may be performed concurrently with a hysterectomy for uterine conditions, such as uterine cancer or large fibroids, to prevent future adnexal issues. Additionally, it may be part of a more extensive cancer debulking surgery where other organs are involved. The determination of its primary or secondary role is based on the underlying diagnosis, the extent of disease, and the overall treatment strategy.

\section*{Relevant Surgical Anatomy}
The surgical anatomy relevant to salpingo-oophorectomy includes the ovaries, fallopian tubes, and surrounding structures within the female pelvis. The ovaries are almond-shaped organs located bilaterally, responsible for hormone production and oocyte release. The fallopian tubes extend from the uterus to the ovaries, facilitating the transport of ova and sperm. 

The arterial supply to the ovaries is primarily via the ovarian artery, which branches from the abdominal aorta, while venous drainage occurs through the ovarian vein, which drains into the inferior vena cava on the right and the renal vein on the left. The lymphatic drainage follows the arterial supply, with lymph nodes located in the para-aortic and pelvic regions.

Key surgical landmarks include:

\begin{itemize}
  \item The infundibulopelvic ligament contains the ovarian vessels; the utero-ovarian ligament and mesosalpinx connect the adnexa to the uterus and broad ligament.
  \item The ureter, iliac vessels, bladder, and bowel are important adjacent structures during pelvic dissection.
\end{itemize}

Understanding these anatomical relationships is crucial for avoiding complications during surgery, such as injury to adjacent structures like the ureters, bladder, and major blood vessels.

\section*{Preoperative Workup \& Preparation}
Thorough preoperative preparation is essential to ensure patient safety and optimize surgical outcomes. This typically includes:

\begin{itemize}
  \item A comprehensive history and physical examination.
  \item Routine laboratory tests such as a complete blood count, electrolyte panel, and coagulation studies.
  \item Specific imaging studies like pelvic ultrasound, CT scan, or MRI to characterize the adnexal mass or assess for malignancy.
  \item Tumor markers (e.g., CA-125, HE4) may be drawn if there is suspicion of cancer.
  \item Anesthesia clearance, which may involve cardiac or pulmonary evaluations if significant comorbidities exist.
  \item NPO (nil per os) instructions for a specified period, usually 6-8 hours before surgery.
  \item Adjustments or temporary discontinuation of medications, particularly anticoagulants.
  \item Administration of preoperative antibiotics intravenously shortly before the incision to reduce the risk of infection.
  \item Detailed informed consent covering risks, benefits, alternatives, and potential implications of the surgery, especially regarding fertility and surgical menopause if bilateral oophorectomy is planned.
\end{itemize}

\textbf{Situations requiring treatment, optimization, or a different approach:}
\begin{itemize}
  \item Uncorrectable severe coagulopathy or bleeding disorder.
  \item Severe, unstable medical conditions that preclude safe general anesthesia or surgery, such as acute myocardial infarction, uncontrolled sepsis, or severe respiratory failure.
  \item Patient refusal after thorough counseling.
\end{itemize}

\textbf{Relative Contraindications and Precautions:}
\begin{itemize}
  \item Severe cardiopulmonary disease, significantly increasing anesthetic and surgical risk, requiring extensive preoperative optimization.
  \item Extensive intra-abdominal adhesions from previous surgeries, endometriosis, or inflammatory conditions, complicating laparoscopic access and dissection.
  \item Active systemic infection or localized pelvic infection (e.g., unruptured tubo-ovarian abscess) where delaying surgery might be beneficial.
  \item Desire for future fertility requires explicit counseling because bilateral removal causes permanent infertility and loss of ovarian hormone production; ovarian-sparing alternatives should be considered when safe.
  \item Morbid obesity, which can present technical challenges for both laparoscopic access and surgical exposure.
  \item Young age for bilateral oophorectomy, due to long-term health implications of premature surgical menopause, necessitating careful counseling and consideration of hormone replacement therapy.
\end{itemize}

\section*{Surgical Technique \& Procedure}
Salpingo-oophorectomy is typically performed under general anesthesia. The patient is positioned supine, often in a Trendelenburg position to facilitate pelvic access, and may be placed in dorsal lithotomy stirrups if uterine manipulation is required. The surgical approach can be either minimally invasive (laparoscopic or robotic) or open (laparotomy).

For a laparoscopic approach, the following steps are generally followed:

\begin{itemize}
  \item A small incision is made at the umbilicus to insert a trocar and establish pneumoperitoneum with carbon dioxide gas.
  \item Additional 2-3 small incisions are made in the lower abdomen for accessory trocars, allowing the insertion of surgical instruments and a camera.
  \item A thorough pelvic survey is conducted to identify the uterus, ovaries, and fallopian tubes.
  \item The infundibulopelvic ligament, containing the ovarian artery and vein, is carefully identified, coagulated, and divided using an energy device (e.g., bipolar cautery, ultrasonic shears).
  \item The utero-ovarian ligament and mesosalpinx are similarly coagulated and divided, detaching the ovary and fallopian tube from the uterus.
  \item The excised adnexa are placed into an endoscopic retrieval bag and removed through one of the port sites.
  \item Hemostasis is confirmed, and the abdomen is deflated before closing the incisions.
\end{itemize}

In an open approach, the procedure involves:

\begin{itemize}
  \item Making a Pfannenstiel (bikini line) or vertical midline incision in the lower abdomen.
  \item Incising the abdominal wall layers and entering the peritoneal cavity.
  \item Direct visualization and mobilization of the adnexa.
  \item Clamping, ligating with sutures, and dividing the infundibulopelvic ligament and utero-ovarian ligament.
  \item Removing the specimen and achieving meticulous hemostasis.
  \item Irrigating the surgical site and closing the incisions in layers.
\end{itemize}

Drains are rarely used for uncomplicated salpingo-oophorectomy.

\section*{Complications \& Risks}
As with any surgical procedure, salpingo-oophorectomy carries potential risks, which can be broadly categorized as general surgical risks and procedure-specific risks.

\textbf{General Surgical Risks:}
\begin{itemize}
  \item \textbf{Bleeding:} Hemorrhage during or after surgery, potentially requiring blood transfusion or reoperation, and hematoma formation at the surgical site.
  \item \textbf{Infection:} Wound infection, pelvic infection, or abscess formation, which may require antibiotics or drainage.
  \item \textbf{Anesthesia Complications:} Adverse reactions to anesthetic agents, respiratory or cardiac complications, nausea, vomiting, or aspiration.
\end{itemize}

\textbf{Procedure-Specific Risks:}
\begin{itemize}
  \item \textbf{Injury to Adjacent Organs:} Accidental damage to the ureter, bladder, bowel, or major blood vessels during dissection.
  \item \textbf{Nerve Injury:} Damage to peripheral nerves from positioning or surgical manipulation, leading to numbness or weakness.
  \item \textbf{Adhesion Formation:} Development of scar tissue within the abdomen, potentially causing chronic pain or bowel obstruction in the future.
  \item \textbf{Ovarian Remnant Syndrome:} Incomplete removal of ovarian tissue, which can continue to produce hormones or develop cysts, requiring further surgery.
  \item \textbf{Port Site Hernia:} Herniation of abdominal contents through a trocar incision site, requiring surgical repair.
  \item \textbf{Conversion to Open Surgery:} The need to switch from a laparoscopic or robotic approach to an open laparotomy due to unforeseen complications or extensive adhesions.
  \item \textbf{Surgical Menopause:} Immediate onset of menopause symptoms if both ovaries are removed, with long-term risks of osteoporosis and cardiovascular disease.
  \item \textbf{Failure to Achieve Desired Outcome:} Recurrence of endometriosis symptoms or incomplete cancer staging/debulking.
\end{itemize}

\section*{Postoperative Care \& Recovery}
Following salpingo-oophorectomy, patients are monitored in the Post-Anesthesia Care Unit (PACU) for vital signs, pain levels, and recovery from anesthesia. Pain management is initiated with intravenous or oral analgesics, and anti-nausea medications are provided as needed.

During the first 24-48 hours, patients are encouraged to ambulate early to prevent complications like deep vein thrombosis. Diet is gradually advanced from clear liquids to solid foods as tolerated. Pain control is transitioned to oral medications, and wound care instructions are provided. Most patients undergoing laparoscopic salpingo-oophorectomy are discharged the same day or within 24 hours. 

Recovery at home typically involves limiting strenuous activities, heavy lifting, douching, tampon use, and sexual intercourse for 4-6 weeks. If bilateral oophorectomy was performed, discussions about managing surgical menopause symptoms and the potential initiation of hormone replacement therapy (HRT) will begin. Full recovery and return to normal activities usually occur within 2-4 weeks for laparoscopic surgery and 4-6 weeks for open surgery.

The most critical result following a salpingo-oophorectomy is the pathology report from the excised ovarian and fallopian tube tissue. This report provides a definitive diagnosis, confirming whether the condition was benign (e.g., endometrioma, dermoid cyst, functional cyst) or malignant (e.g., epithelial ovarian cancer, fallopian tube cancer, borderline tumor). 

For benign conditions, a clear pathology report typically confirms the successful removal of the diseased tissue and guides any further non-surgical management. If malignancy is identified, the pathology report will detail the tumor type, grade, and extent of disease, which is crucial for cancer staging and determining the need for adjuvant therapies such as chemotherapy or radiation. The discharge summary will also provide a comprehensive overview of the surgical findings, any intraoperative complications, and detailed postoperative instructions for the patient's recovery and follow-up care.

A normal, successful postoperative course following salpingo-oophorectomy is characterized by the resolution of preoperative symptoms, such as pelvic pain or abnormal bleeding. Patients typically experience a gradual decrease in pain levels, with most reporting significant relief within the first few days post-surgery. 

The timeline for healing varies based on the surgical approach, with laparoscopic patients often returning to normal activities within 2-4 weeks, while those undergoing open surgery may take 4-6 weeks. Monitoring for complications is essential, and patients are advised to report any signs of infection, excessive bleeding, or unusual pain promptly.

Postoperative follow-up is crucial for monitoring recovery and addressing any long-term implications. 

\begin{itemize}
  \item \textbf{Postoperative Visits:} A first follow-up visit is typically scheduled 1-2 weeks after surgery for a wound check and to discuss initial recovery. A second visit at 4-6 weeks assesses full recovery, reviews the final pathology report, and addresses any ongoing concerns.
  \item \textbf{Suture/Staple Removal:} If non-absorbable sutures or staples were used for skin closure, they will be removed at the first follow-up visit.
  \item \textbf{Imaging and Laboratory Tests:} For patients with a diagnosis of malignancy, regular surveillance imaging (e.g., CT scans, PET scans) and monitoring of tumor markers (e.g., CA-125) will be initiated according to established cancer guidelines.
  \item \textbf{Rehabilitation/Physical Therapy:} Rarely needed for uncomplicated salpingo-oophorectomy, but may be recommended after extensive surgery or if specific functional deficits arise.
  \item \textbf{Activity Resumption:} Patients are advised on a gradual return to normal activities, with full activity, including exercise and sexual intercourse, typically permitted by 4-6 weeks post-surgery.
  \item \textbf{Warning Signs:} Patients are educated on warning signs that warrant immediate medical attention, such as fever, severe or worsening abdominal pain, heavy vaginal bleeding, foul-smelling wound discharge, or persistent difficulty with urination or bowel movements.
  \item \textbf{Management of Surgical Menopause:} For patients who underwent bilateral oophorectomy, ongoing discussions regarding hormone replacement therapy (HRT) to mitigate menopausal symptoms and long-term health risks (e.g., osteoporosis screening, cardiovascular health monitoring) are essential.
\end{itemize}

\section*{Outcomes \& Prognosis}
Salpingo-oophorectomy is one of the most commonly performed gynecologic surgical procedures, often undertaken in conjunction with a hysterectomy. The incidence of salpingo-oophorectomy increases with age, reflecting the higher prevalence of both benign and malignant adnexal pathologies in older women. 

While benign indications such as ovarian cysts, endometriosis, and tubo-ovarian abscesses are more frequent in women of reproductive age, the incidence of ovarian and fallopian tube cancers peaks in post-menopausal women. Risk-reducing salpingo-oophorectomy, however, is typically performed in pre-menopausal women with genetic predispositions, often between the ages of 35 and 45, to maximize cancer prevention benefits before the typical onset of these cancers. 

The overall mortality rate for elective salpingo-oophorectomy is very low, particularly with minimally invasive approaches, while morbidity rates vary depending on the complexity of the case, the underlying indication (benign vs. malignant), and patient comorbidities.

The outcomes and prognosis following salpingo-oophorectomy are generally favorable, though they vary significantly based on the underlying indication. For benign conditions, the typical success rate for symptom resolution (e.g., pain relief, resolution of mass effect) is high, and recurrence of the specific adnexal pathology is rare after complete removal. 

In cases of endometriosis, while the removed endometrioma will not recur, symptoms may persist or recur if other endometriotic implants remain elsewhere in the pelvis. For malignancies, the prognosis is highly dependent on the stage and grade of the cancer at diagnosis, the completeness of surgical removal, and the response to adjuvant therapies. Early-stage ovarian or fallopian tube cancer often has a good prognosis, while advanced stages carry a higher risk of recurrence and lower survival rates. 

Risk-reducing salpingo-oophorectomy substantially lowers, but does not eliminate, tubo-ovarian cancer risk in appropriately selected high-risk individuals. Bilateral premenopausal oophorectomy causes immediate surgical menopause and requires counseling about vasomotor symptoms, bone, cardiovascular, cognitive, sexual, and fertility consequences; hormone therapy may be appropriate for some patients after individualized risk assessment.

\section*{References}
\begin{enumerate}
  \item MedlinePlus. Salpingo-oophorectomy. U.S. National Library of Medicine; 2024. Available from: \url{https://medlineplus.gov/ency/article/002930.htm}.
  
  \item Berek JS, Hacker NF. Berek \& Hacker's Gynecologic Oncology. 6th ed. Wolters Kluwer; 2015.
  
  \item American College of Obstetricians and Gynecologists (ACOG). Practice Bulletin No. 189: Management of Adnexal Masses. Obstet Gynecol. 2018 Mar;131(3):e93-e104.
  
  \item National Comprehensive Cancer Network (NCCN). NCCN Clinical Practice Guidelines in Oncology: Ovarian Cancer. Version 1.2024. Available from: \url{https://www.nccn.org/professionals/physician_gls/pdf/ovarian.pdf}.
\end{enumerate}

\clearpage
